\documentclass{article}
\title{Nyasha}
\begin{document}
\section*{Nyasha, Oh Starving!}
\textit{``They do not like my language , my English, because it is authentic and my Shona, because it is not!'' }
\paragraph{}

Nyasha is the character most affected by the titular `nervous conditions'. For Nyasha food is symbolic of so much; locale since she has "tasted" what it feels like to not be in Zimbabwe. This is seen when she returns from England. \textit{``Nyasha tucked into the vegetables with the rest of us. When my mother offered her the sour milk she had asked for, she became very morose. She refused to eat'' }. Nyasha wishes to not be in the physical space that she is and so she rejects the food that is representative of that. Food to Nyasha is also a cudgel her father uses to control her. 
\paragraph{}
\textit{``I'm not hungry,' Nyasha explained.}\\

\textit{"You will eat that food,' commanded the man. "Your mother and I are not killing ourselves working just for you to waste your time playing with boys and then come back and turn up your nose at what we offer. Sit and eat that food. I am telling you. Eat it!' Nyasha took a few mouthfuls.}\\

\textit{'She has had enough, Baba,' Maiguru said, but Babamukuru was adamant. He was very upset.}\\

\textit{'She must eat her food, all of it. She is always doing this, challenging me. I am her father. If she doesn't want to do what I say, I shall stop providing for her - fees, clothes, food, everything.'}\\


\textit{'Nyasha, eat your food,' her mother advised.}
\textit{'Christ!' Nyasha breathed and with a shrug picked up her fork and began to eat, slowly at first, then gobbling the food down without a break. The atmosphere lightened with every mouthful she took.''}
\paragraph{}

While an entire essay could be written on this scene; for this essay, it lays bare how ``the man'' sees not eating the food he provides as a challenge to his authority. In this moment he is not Babamukuru, Nyasha's father, but instead the man or even just men. After Nyasha's last attempt at outward rebellion she almost lost her life. So here she cannot rebel outwardly and instead waits later to vomit it up. The path Nyasha cuts is the most damaging, yet it allows her to retain her dignity. This is solidified when Nyasha later writes to Tambu and says ``[Nyasha] had embarked on a diet "to discipline my body and occupy my mind" '' Nyashas is characterized as a guerilla fighter, and sees her anorexia as training for a final conflict. Continuing to inspect food as means for control we know Nyasha has no way to make food on her own. She lacks the skills to grow vegetables that Tambu has or Lucias apparent ability to conjure it. ``Can you cook books'' is what Tambus father says to her but in the end this applies far better to Nyasha, she is incredibly educated and in the peak of her anorexia she appears to be attempting to live off books. For all Nyashas education, she cannot provide food for herself, so instead she develops anorexia as means to control her destiny.


\end{document}